\documentclass[11pt,a4paper]{article}
\usepackage[margin=2.6cm]{geometry}
\usepackage{amsmath,amssymb,amsthm}
\usepackage{booktabs}
\usepackage[hidelinks]{hyperref}
\newtheorem{proposition}{Proposition}
\title{No four coplanar in the cube: a witness for $a(7)\ge18$ and the state of the upper bound}
\author{Aleksei Kudriashov (Alex Komang)\\ \small Nusa Dua Studio, Nusa Dua, Bali, Indonesia\\ \small ORCID: 0009-0006-8885-5338}
\date{20 August 2026 --- version 1.1}
\begin{document}\maketitle
\begin{abstract}
Let $a(n)$ be the largest number of points of $\{0,\dots,n-1\}^3$ no four of which are coplanar, OEIS
A280537.  The entry records $5,8,10,13,16,18,20$ for $n=2,\dots,8$ and states that the terms up to $a(6)$ were
found by exhaustive search while $a(7)$ and $a(8)$ rest on numerical evidence.  We reproduce $a(2),\dots,a(5)$
by three independent methods --- $a(5)=13$ with an unsatisfiability certificate verified by an external
checker --- exhibit witnesses for $a(6)\ge16$ and $a(7)\ge18$, which the entry did not carry, and report the
state of the remaining question, the impossibility of $19$ points at $n=7$.

We also record how we nearly published the opposite.  Having closed $124$ of $128$ regions of a case split at
$M=18$ without finding a configuration, we leaned towards $a(7)=17$ and towards announcing that the published
value is wrong.  The four unclosed regions were precisely the ones the symmetry pruning drives solutions into,
and the configuration was in them.  It was found by an independent search and verified here three ways.  Nothing
in our computations was incorrect; what was incorrect was reading the fraction of closed regions as an
approximation to a conclusion.
\end{abstract}

\section{Reformulation}
Four points are coplanar exactly when they lie in a common plane, and a plane carrying at most three lattice
points forbids nothing.  Hence
\[S\subseteq[n]^3\text{ is admissible}\iff |S\cap P|\le3\ \text{ for every plane }P\text{ with}\ge4\text{ lattice points,}\]
a propositional formula on $n^3$ variables.  The layer $z=k$ is itself such a plane, so $a(n)\le 3n$; the
published values satisfy $5\le6$, $8\le9$, $10\le12$, $13\le15$, $16\le18$, $18\le21$, $20\le24$.

\section{What is established here}
\begin{center}\begin{tabular}{rrl}\toprule
$n$ & value & how\\\midrule
2 & $a=5$ & three independent methods\\
3 & $a=8$ & three independent methods; solution count matches the published one\\
4 & $a=10$ & three methods; \textbf{232 classes}, matching the published count\\
5 & $a=13$ & three methods; unsatisfiability at $14$ carries a certificate verified by \texttt{drat-trim}\\
6 & $a\ge16$ & four independent witnesses, no two equivalent under the cube group\\
7 & $a\ge18$ & witness, verified three ways here and four ways independently\\\bottomrule
\end{tabular}\end{center}
The strongest of these checks is not that the maxima agree but that the \emph{counts} do.  Ed Pegg reported a
unique $8$-point configuration at $n=3$, $232$ configurations of $10$ points at $n=4$, and $38$ ways to choose
$13$ points at $n=5$.  Our enumerator gives
\begin{center}\begin{tabular}{rrrr}\toprule
$n$ & size & labelled configurations & classes under the cube group\\\midrule
3 & 8 & 16 & \textbf{1} (published: unique)\\
4 & 10 & 10960 & \textbf{232} (published: 232)\\
5 & 13 & 1768 & \textbf{38} (published: 38)\\\bottomrule
\end{tabular}\end{center}
with the stabiliser checksum $\sum 48/|\mathrm{Stab}|$ equal to the labelled total in each row --- at $n=4$,
$225\cdot48+6\cdot24+1\cdot16=10960$.  Matching a maximum can happen by accident; matching a count three times
cannot.  We stress that the enumerator producing these numbers has no notion of a plane at all: it tests
coplanarity by an integer determinant on each quadruple, so agreement with our plane-based encoding is
agreement between two constructions that share nothing but the problem.

\section{Three values proved twice, by paths sharing no code}
For $n\le5$ the values do not depend on whether the propositional machinery of this note is sound at all.  A
second route establishes them without a solver, without an encoding and without any notion of a plane:
coplanarity is tested by an integer determinant on each quadruple, and the search is split over orbits of
\emph{triples} of layer profiles under the cube group.
\begin{center}\begin{tabular}{rrrrl}\toprule
$n$ & $M$ & orbits & weight sum vs.\ triples & configurations found\\\midrule
3 & 9 & 1 & $1/1$ & 0 $\Rightarrow a(3)\le8$\\
4 & 11 & 4 & $64/64$ & 0 $\Rightarrow a(4)\le10$\\
5 & 14 & 10 & $125/125$ & 0 $\Rightarrow a(5)\le13$\\\bottomrule
\end{tabular}\end{center}
The weight sum is checked against the total number of triples in every row, so no orbit is missing; omitting one
would void the argument.  With the witnesses already exhibited, this gives $a(3)=8$, $a(4)=10$, $a(5)=13$ twice
over, by constructions that share nothing but the problem statement.

The distinction from the count agreement above is worth stating.  Matching counts says the two enumerators see
the same problem.  An independent proof says the result does not depend on our propositional machinery being
correct at all.  The boundary is sharp: at $n=6$ this route was measured and does not reach (some $70$ hours on
four cores), and at $n=7$ neither route has an upper bound.  The double proof covers $n\le5$ and not one value
further.

\section{What is \emph{not} established: there are no upper bounds here}
This deserves its own paragraph, because a reader who has seen the companion note on the no-three-in-line
problem --- where the upper bounds are the substance --- will otherwise supply one by habit.

For the present problem we have \textbf{no completed upper bound above $n=5$}.  At $n=5$ the impossibility of
$14$ points is established and carries a certificate.  At $n=6$ we have four witnesses for $16$ and a case
split at $M=17$ that is not finished; the published $a(6)=16$ rests on someone else's exhaustive search, not
on ours.  At $n=7$ we have six witnesses for $18$ and a case split at $M=19$ in progress.  Every entry of the
table above that is written $a\ge$ is a lower bound and nothing more.

We say this explicitly because the tempting summary --- ``we settled $a(7)$'' --- would be false in the half
that is hard.  Exhibiting a configuration is cheap and self-verifying; showing that none is larger is neither.

A second admission belongs here.  At $n=6$ the published $a(6)=16$ is the one value in this sequence that was
proved before us, by exhaustive search, and it is therefore our calibration: reproducing it is what licenses
anything we say about $n=7$.  We reproduce it by one route only.  A second, propositional-logic-free route was
attempted --- a direct enumeration testing coplanarity by determinants, with no notion of a plane and no solver
--- and abandoned after its cost was measured rather than waited out: a calibrated tree-size estimator puts it
at some $28$ hours per orbit representative and $70$ hours in total on four cores, and that figure is a lower
bound, since the ratio of mean to median across probes is $31$ and a heavy tail makes a finite sample
systematically optimistic.  So the calibration is single-path.  We prefer saying so to the appearance of double
confirmation.

\section{The near-miss, in full}
We record this because the mechanism is more transferable than the result.

The decisive question was split into regions by the contents of the first columns; a column cannot hold four
points, since four points of a column are collinear and hence coplanar, so subsets of size at most three
exhaust it.  At $M=18$ we closed $124$ of $128$ regions with no configuration found, and reported the fraction
as though it were progress towards a conclusion.

It was not.  The pruning we use keeps, from each orbit of the cube group, the lexicographically least member ---
the one with zeros as early as possible.  Solutions are therefore driven systematically into the least
constrained regions, which are exactly the ones that close last.  The unclosed regions were not a remainder;
they were the whole question.  We had written this out for each other two hours earlier and still read the
fraction the other way.

An independent search then produced three $18$-point configurations.  We verified one here by three separate
computations: all $\binom{18}{4}=3060$ quadruples with the $3\times3$ determinant of differences, all $3060$
with the $4\times4$ determinant of homogeneous coordinates, and membership in our list of $346743$ rich planes
at $n=7$, itself verified complete against $6650464$ non-collinear triples.  All three give zero.

What saved us was not caution.  It was that neither of us ever wrote ``does not exist'' where we meant ``have
not found''.

\section{The remaining question}
Whether $19$ points fit at $n=7$ --- equivalently, whether $a(7)\le18$ --- is open at the time of writing.  An
exhaustive case split is running; $402$ of its $448$ regions are closed and none is satisfiable.  We deliberately
do not offer that fraction as evidence: the unclosed regions are the least constrained ones, which is precisely
where the symmetry pruning drives solutions, and reading such a fraction as progress towards a conclusion is the
error described in the previous section.

We can, however, say what it costs.  The regions still open subdivide into pieces of which roughly half are
hard, not one in sixty-four as at earlier levels, and the hard ones run beyond ten minutes each; the projection
is some eight hours on the forty-four cores available, and the tail of such distributions is not to be trusted
downwards.  A value that has stood unproved for nine years is unlikely to fall in an evening.  Should the split
close, this note will be revised under the same concept DOI.

\section*{AI/tool-use disclosure}
The searches, the programs and the text were produced by autonomous AI agents (Anthropic Claude) under the
direction of the author of record, who posed the question, chose what to compute, decided what to claim and
takes responsibility for the content.  Verification protocol, programs, journals and witnesses:
\url{https://github.com/iwasborninbali/saturation}.
\end{document}
